%% Withdrawal Reactions from Social Networks
\input template-en.tex

\title{Withdrawal Reactions from}
\title{Social Networks}
\date{April 18, 2017}

In some cyberpunk-themed works, almost everyone is equipped with a brain-computer interface, communicating via cables plugged into the back of their necks. Although this seems extremely science-fictional, when I see crowds on the subway snickering at their phones, I feel that this neck-plugged cyberpunk has already appeared in reality through some degraded form. This is the unique digital survival mode of our cyber-era Social Networking Services.

Over the last two days, I conducted an experiment: completely disconnecting from social networks. Of course, the experiment failed; after holding out for less than 60 hours, I couldn't help but open Weibo. However, there are things worth reflecting on in this process.

A comprehensive, blind-spot-free block of social networks did achieve some results. On the first day of the experiment, I intentionally filled my schedule so that I wouldn't have any spare time to think about social networks, and I went to bed early at night. This progress naturally made me very happy, thinking that the charm of social networks was just that. But subsequent developments made me realize I was still ``too young, too simple.''

Starting from the second day, I felt significant discomfort. At first, it was just a slight uneasiness, similar to what many articles discussing social networks mention—the fear of ``missing out'' (FOMO) on something. But it turns out this experience isn't that simple. I told myself rationally countless times that what was happening on social networks at that moment was the same as usual—trivial matters. However, this reason was ineffective; the craving for social networks didn't come from reason, but felt more like a longing from the subconscious, from instinct.

By nightfall, things became more serious. Although I didn't wake up late and my routine was fairly regular, I couldn't fall asleep. A desire to scroll through Weibo completely occupied my mind. Though I eventually restrained myself, hours had passed before I fell asleep.

Finally, on the third day, due to lack of rest, my self-control completely collapsed. After clearing the relevant entries in the hosts file, I started scrolling through Weibo again. The action of breaking away from social networks was thus declared a failure.

It seems the significance of detaching from social networks is not obvious, and such behavior even appears somewhat ridiculous, a bit like the ``abstinence'' advocated by Baidu's Jie Se (Abstinence) Bar. However, in a significant sense, social networks have indeed seriously affected my life. It's hard to count exactly how much time I spend on social networks every day, but one thing is certain: the current situation is abnormal. Usually, scrolling through social networks is considered an entertainment activity, but it cannot bring relaxation like singing or exercise. On the contrary, in many cases, it increases feelings of tension. Therefore, reducing the use of social networks doesn't have any ``successology'' significance; in other words, even if I truly and completely stopped using social networks, I wouldn't necessarily use that time for studying or practice. But it can bring a change in lifestyle, moving away from obsessive-compulsive morbidity and facing one's inner self directly.

I have no experience with addictions like tobacco or alcohol, but when I look back at the various reactions I had during those two or three days of being away from social networks, I feel an inexplicable similarity between the two. Many studies have also confirmed my suspicion; according to an experiment at the University of Chicago, Twitter is harder to resist than tobacco and alcohol.

\bye